% !TEX root = ../main.tex

\chapter{Design Specification}

\section{Stakeholders and Workflow}

The primary stakeholder is a researcher who needs to triage new papers, inspect a paper's contribution, ask evidence-seeking questions, revisit related work, and refine future discovery.
The supplied materials also imply secondary stakeholders: project teams that may share papers, system operators who manage model and storage costs, and paper authors whose work must not be misrepresented.
Because raw interview artifacts were not supplied, this chapter does not report interview counts, participant quotations, or prevalence estimates.
Instead, it treats the workflow encoded in the story map and prototype as a design scenario to be evaluated.

The intended workflow begins with a digest.
Each recommended paper includes metadata, a concise summary, and a recommendation reason.
Opening the paper reveals a structured summary and evidence controls.
A question opens a scoped retrieval-and-generation interaction.
The user may then inspect citation relations or modify interest settings.
Events generated by these actions can inform later recommendations, but only according to the event semantics defined in Chapter 3.

\section{Functional Requirements}

Table~\ref{tab:functional-requirements} condenses the feature plan into research-facing requirements.
The requirements are grouped by the invariant they protect rather than by implementation library.

\begin{longtable}{P{1.0cm}P{4.7cm}P{5.0cm}P{2.8cm}}
\caption{Functional requirements mapped to the thesis challenges.}
\label{tab:functional-requirements}\\
\toprule
ID & Requirement & Acceptance evidence & Traceability \\
\midrule
\endfirsthead
\toprule
ID & Requirement & Acceptance evidence & Traceability \\
\midrule
\endhead
R1.1 & Assign a stable logical identity and an explicit revision identity to every ingested paper. & Version fixtures show that two revisions remain distinguishable while belonging to one logical paper. & TC1--E1 \\
R1.2 & Propagate revision and transformation metadata to sections, chunks, summaries, embeddings, and answer citations. & A provenance query reconstructs every dependency for a displayed claim. & TC1--E1 \\
R1.3 & Label parsing quality and retain a defined degraded mode. & Multi-layout fixtures produce full, partial, or abstract-only states without silent success. & TC1--E1 \\
R2.1 & Scope retrieval by paper, selected library, or declared collection. & Retrieval traces contain only documents allowed by the selected scope. & TC2--E2 \\
R2.2 & Represent generated answers as claims linked to source spans. & Every externally checkable claim contains resolvable evidence or an explicit unsupported state. & TC2--E2 \\
R2.3 & Distinguish identifier, availability, consistency, and coverage checks. & Failure fixtures trigger the intended verification label. & TC2--E2 \\
R2.4 & Abstain when available evidence does not support an answer. & Unanswerable questions are evaluated separately from answerable questions. & TC2--E2 \\
R3.1 & Record explicit and implicit events with documented semantics and deduplication identifiers. & Replayed event batches are idempotent and event provenance is inspectable. & TC3--E3 \\
R3.2 & Maintain one bounded, time-decayed normalized behavior-v1 vector; keep separable multi-timescale state as planned behavior-v2. & Synthetic event replays yield predictable bounded and decayed behavior-v1 updates. & TC3--E3 \\
R3.3 & Explain recommendation reasons using inspectable evidence from the interest state. & Reasons map to declared topics, authors, papers, or events. & TC3--E3 \\
R3.4 & Permit users to edit, suppress, or delete inferred interests and behavioral history. & Correction tests confirm that the visible and ranking states change consistently. & TC3--E3 \\
R4.1 & Preserve essential reading state during intermittent connectivity. & Offline and reconnection scenarios retain consistent local state. & TC4--E4 \\
R4.2 & Bound background work, push frequency, and model cost. & Configuration and logs demonstrate rate, frequency, and budget limits. & TC4--E4 \\
R4.3 & Expose degraded states and verification actions in the mobile UI. & Usability tasks measure whether users can identify evidence and recover from failures. & TC4--E4 \\
\bottomrule
\end{longtable}

\section{Non-Functional and Ethical Requirements}

\subsection{Provenance and Reproducibility}

Every claim-bearing artifact must record enough configuration to be reproduced: source document identity and revision, parser version, chunking configuration, embedding model, retrieval parameters, generator and prompt version, and verification policy.
The requirement does not imply that all third-party models are deterministic.
It requires the system to preserve the conditions under which an output was produced.

\subsection{Privacy and Data Minimization}

Behavioral events are potentially sensitive because they reveal research topics, reading habits, and professional activity.
The event schema should therefore collect only fields required by the declared update model.
Users must be able to inspect and remove events and inferred interests.
Retention, export, deletion, and synchronization policies must be explicit.
The system must not interpret a missing click as a stable negative preference without a stated exposure condition.

\subsection{Uncertainty and User Control}

Human--AI interaction guidelines emphasize communicating capability, supporting correction, and responding to feedback~\cite{amershi2019guidelines}.
For MNEME, uncertainty is not represented by a single confidence number.
The interface should instead expose operational states: evidence retrieved, source span available, consistency check passed or failed, answer incomplete, or no answer supported.
Users should be able to open the source, change retrieval scope, reformulate a question, and correct the interest state.

\subsection{Latency, Reliability, and Cost}

Latency targets must be measured per operation rather than summarized by one end-to-end average.
Digest loading, paper detail retrieval, local navigation, first-token latency, complete-answer latency, graph expansion, and synchronization have different user consequences.
Reliability similarly requires error classes: network failure, provider failure, parse failure, retrieval-empty, verification failure, and local synchronization conflict.
Cost boundaries should include a daily model budget and cached results keyed by content and configuration hashes.

\section{Interaction Design}

\subsection{Digest and Paper Detail}

The interactive prototype presents the digest as the home screen.
A paper card exposes its title, authors, a concise description, a match indicator, a recommendation reason, and an open action.
This organization is intended to support rapid triage.
The paper-detail view then separates summary content from source-verification actions so that the user can move from overview to evidence.
These statements describe the prototype and do not establish production behavior.

\subsection{Cited Question Answering}

The Q\&A view is intended to display an answer together with source chips or evidence spans.
The interaction must not label an answer ``verified'' merely because a paper identifier is valid.
The user-visible status should correspond to the verification layers defined in Chapter 3.
Multi-turn interaction must also preserve the selected scope and make scope changes visible.
Recent multi-turn RAG evaluation shows that follow-up, unanswerable, and non-standalone questions create distinct challenges~\cite{katsis2025mtrag}; therefore, follow-up convenience cannot replace explicit evaluation.

\subsection{Knowledge Graph and Interest Controls}

The graph view presents related or citing papers as nodes.
The interest view permits topic and preference adjustments.
These screens support two different forms of continuity: the graph externalizes relations among papers, while the interest view externalizes part of the system's user model.
Both require clear actions.
A visual node without a visible selected state or primary action can appear decorative rather than interactive.
An interest label without provenance or correction controls can appear authoritative even when inferred from weak evidence.

\section{Formative Prototype Evidence}

The final usability presentation reports that five participants completed five prototype tasks.
It reports success counts of 5/5 for opening a recommended paper, 5/5 for identifying evidence, 5/5 for cited Q\&A, 2/5 for opening a related paper from the graph, and 4/5 for changing preferences.
The presentation interprets graph discoverability as the primary drop-off and proposes an instruction, a selected-node state, and a persistent ``Open paper'' action.

These counts are classified as \observed{}/\limited{}, not \measured{} final evidence.
The supplied folder does not include the participant-by-task table, timing and tap records, coded observations, consent documentation, or locally accessible recordings.
The companion draft presentation explicitly contains blank fields for those artifacts.
Consequently, the counts can motivate a design iteration, but they cannot establish a general usability effect or validate the redesign.

\section{Mobile Acceptance and Evidence Boundaries}

The mobile requirements become testable when each one names a reader-visible result under a defined condition.
Table~\ref{tab:mobile-acceptance} extends the existing R4 requirements with client-side acceptance conditions.
It specifies what the Android client must expose; Chapter~4 reports the corresponding E4 measurements.

\begin{table}[htbp]
\centering
\caption{Client-side acceptance conditions for bounded mobile integration.}
\label{tab:mobile-acceptance}
\begin{tabular}{P{2.4cm}P{6.2cm}P{5.3cm}}
\toprule
Requirement link & Observable acceptance condition & E4 evidence \\
\midrule
R4.1: reading continuity &
A stored backend briefing remains available when refresh fails, carries a cached-content label, and is never substituted for live content without disclosure.
When no valid stored briefing exists, the client presents a retryable error. &
Controlled live, cached, offline, server-failure, and invalid-cache trials. \\
\addlinespace
R4.1 and R4.3: asynchronous work &
Paper preparation exposes loading, processing, ready, and error states.
Polling terminates at a completed or failed job and does not leave the interface in an unexplained waiting state. &
Controlled three-stage polling and startup/resume measurements. \\
\addlinespace
R4.3: graph interaction &
The graph declares ready, fallback, or empty state.
A selected node remains visible in the native interface and provides an Open Paper action that returns to the same graph selection. &
Graph-state trials at the declared client bound and graph/navigation software tests. \\
\addlinespace
R3.1 and R4.1: event return &
Visible paper interactions enter a durable queue with stable identifiers.
Retryable failures preserve pending work, and replayed identifiers are acknowledged as duplicates. &
Controlled queue/retry/replay trials and a separate live-backend trace. \\
\bottomrule
\end{tabular}
\end{table}

These conditions support implementation and systems acceptance.
They do not establish whether readers understand every status label or whether the graph redesign improves task completion.
No participant retest accompanied the Android measurements, so those user-facing effects remain unestablished.

The implemented event payload also narrows the privacy boundary.
Paper impressions, paper opens, and submitted paper-scoped questions are represented by an event type, relevant identifiers, and a timestamp.
The submitted question text is not copied into the synchronization event.
This data-minimization decision reduces the behavioral content held in the client queue; it is not a complete privacy or security audit.
